% bab2.tex — gabungan BAB_II_A/B/C *_POLISHED.md (## A + ## B + ## C)

\phantomsection
\section*{BAB II --- PEMBAHASAN}
\addcontentsline{toc}{section}{BAB II --- PEMBAHASAN}

\phantomsection
\subsection*{A. Periodisasi Sejarah Perkembangan Hadis}
\addcontentsline{toc}{subsection}{A. Periodisasi Sejarah Perkembangan Hadis}

\phantomsection
\subsubsection*{1. Periodisasi Masa Nabi Muhammad SAW (610--632 M / 13 SH--11 H)}
\addcontentsline{toc}{subsubsection}{1. Periodisasi Masa Nabi Muhammad SAW (610--632 M / 13 SH--11 H)}

Ciri khas penyaluran hadis pada era kenabian terletak pada keragaman jalurnya: hafalan, pengamalan ibadah, \textit{qira'ah}, \textit{sama'}, dan pencatatan perdana berjalan beriringan sebagai satu ekosistem penjagaan.\footnote{T. N. S. Tanjung and Z. Fadhlurrahman, "The role of tadwin hadith in the development of Islamic historiography," \\textit{Istifham: Journal of Islamic Studies}, vol. 3, no. 3, hlm. 188--201, 2025; A. Musadad, "Articulating the prophetic authority: Some notes on the nature of hadith collection in al-Muwatta' and musannaf literature," \\textit{Jurnal Studi Ilmu-Ilmu Al-Qur'an dan Hadis}, vol. 27, no. 1, hlm. 351--373, 2026.} Hadis bukan sekadar ujaran tekstual, melainkan \textit{living sunnah} yang dihidupkan lewat suri teladan Nabi dalam salat, muamalah, dan pergaulan sosial.\footnote{N. M. A. M. Razak et al., "Evolution of hadith textual criticism: From classical to contemporary approaches," \\textit{Global Journal Al-Thaqafah}, hlm. 15--29, 2025; Tanjung, "The Role of Tadwin Hadith.".} Kelisanan yang kukuh pada masyarakat Arab yang \textit{ummi} menjadikan ingatan kolektif sebagai sarana utama, sedangkan tulisan sekadar menyangga memori, bukan kodifikasi sistematis.\footnote{Musadad, "Articulating the Prophetic Authority."; M. A. Mosa, "Synergizing structure and semantics: A knowledge graph-transformer framework for narrator disambiguation in hadith networks," \\textit{Digital Scholarship in the Humanities}, 2025.} Kewibawaan bertumpu pada kehadiran ragawi Nabi selaku sumber hidup, sehingga rantai pembuktian formal belum dituntut.\footnote{Tanjung, "The Role of Tadwin Hadith."; Razak, "Evolution of Hadith Textual Criticism.".} Pewarisan hadis menempuh empat jalur: \textit{sama'} (menyimak lalu menghafal), \textit{qira'ah}/\textit{'ard} (melafalkan kembali di hadapan Nabi untuk dikoreksi), pengajaran interaktif melalui majelis dan khotbah, serta penyaksian partisipatif dalam praktik wudu, salat, dan haji.\footnote{Tanjung, "The Role of Tadwin Hadith."; Razak, "Evolution of Hadith Textual Criticism."; Musadad, "Articulating the Prophetic Authority."; Mosa, "Synergizing Structure and Semantics.".} Corak interaktif-dialogis itu, dipadukan dengan ekologi lisan-tulis yang saling melengkapi, menegaskan bahwa pemisahan kaku antara oral dan tertulis tidaklah memadai.\footnote{Razak, "Evolution of Hadith Textual Criticism."; Tanjung, "The Role of Tadwin Hadith."; Tanjung, "The Role of Tadwin Hadith."; Musadad, "Articulating the Prophetic Authority.".} Fase Nabi karenanya tampil sebagai tahapan formatif yang mutlak, mewariskan memori, praktik, dan catatan yang terserak sebagai bahan mentah bagi seluruh penilaian angkatan berikutnya.\footnote{Razak, "Evolution of Hadith Textual Criticism."; Mosa, "Synergizing Structure and Semantics."; Musadad, "Articulating the Prophetic Authority."; Tanjung, "The Role of Tadwin Hadith.".}

Persoalan penulisan hadis diwarnai tarik-menarik dua klaster riwayat: pelarangan mencatat selain Al-Qur'an berhadapan dengan perkenan tegas untuk mendokumentasikan sabda Nabi.\footnote{Tanjung, "The Role of Tadwin Hadith."; Razak, "Evolution of Hadith Textual Criticism.".} Pelarangan awal dimaknai sebagai \textit{sadd al-dzari'ah}, upaya menghindari bercampurnya wahyu yang masih turun bertahap dengan perkataan Nabi di tengah umat yang baru mengenal aksara.\footnote{Tanjung, "The Role of Tadwin Hadith."; Musadad, "Articulating the Prophetic Authority.".} Seiring menguatnya keberaksaraan Madinah dan kesanggupan membedakan Al-Qur'an dari hadis, Nabi memberikan kelonggaran bersyarat (\textit{rukhsah}), sehingga larangan dan perkenan dipahami sebagai strategi berjenjang, bukan pertentangan.\footnote{Razak, "Evolution of Hadith Textual Criticism."; Tanjung, "The Role of Tadwin Hadith.".} Kompromi ini sejalan dengan penilaian bahwa tradisi tulis perdana memang telah ada, meski mutu keaslian tiap naskah bertingkat dan tidak seluruhnya dapat direkonstruksi secara kukuh.\footnote{Musadad, "Articulating the Prophetic Authority."; Razak, "Evolution of Hadith Textual Criticism.".} Klaim yang meniadakan penulisan hadis pada abad pertama Hijriah karenanya terlalu jauh, sebab cikal dokumentasi tertulis telah tumbuh berdampingan dengan kelisanan.\footnote{Tanjung, "The Role of Tadwin Hadith."; Mosa, "Synergizing Structure and Semantics.".}

Kalangan sahabat terbagi dalam dua peran fungsional yang saling melengkapi: penulis dan penghafal.\footnote{Tanjung, "The Role of Tadwin Hadith."; Musadad, "Articulating the Prophetic Authority.".} Abdullah bin Amr bin al-Ash tampil sebagai penulis utama melalui \textit{Sahifah al-Sadiqah} yang direstui Nabi, bersama Piagam Madinah dan surat-surat diplomatik yang menegaskan keberaksaraan komunitas perdana.\footnote{Tanjung, "The Role of Tadwin Hadith."; Razak, "Evolution of Hadith Textual Criticism.".} Para penghafal seperti Abu Hurairah, Aisyah, Anas bin Malik, dan Ibnu Abbas bersandar pada kekuatan memori dan ketekunan bergaul dengan Nabi, sehingga menjadi titik sentral transmisi.\footnote{Mosa, "Synergizing Structure and Semantics."; Musadad, "Articulating the Prophetic Authority.".} Keragaman logat kesukuan dan keterbatasan ingatan menumbuhkan periwayatan \textit{bi al-ma'na} di samping \textit{bi al-lafdzi} --- kemestian sosiolinguistik yang kelak menuntut kritik matan.\footnote{Razak, "Evolution of Hadith Textual Criticism."; M. R. Rozikin, "Al-fara'id is half of knowledge: A critical hadith analysis and its juridical implications in Islamic inheritance discourse," \\textit{Cogent Arts and Humanities}, vol. 13, no. 1, 2026.} Wafat Nabi pada 11 H/632 M mengakhiri otoritas tunggal yang sentralistis dan membuka era sahabat yang tercerai-berai, sehingga pembuktian dan pewilayahan transmisi menjadi agenda pokok berikutnya.\footnote{Tanjung, "The Role of Tadwin Hadith."; H. H. Liew, "The caliphate will last for thirty years: Polemic and political thought in the afterlife of a prophetic hadith," \\textit{Journal of Islamic Studies}, vol. 36, no. 1, hlm. 38--82, 2025.}

\phantomsection
\subsubsection*{2. Periodisasi Masa Sahabat (11--80 H / 632--700 M)}
\addcontentsline{toc}{subsubsection}{2. Periodisasi Masa Sahabat (11--80 H / 632--700 M)}

Watak menonjol era sahabat ialah penyebaran geografis sepeninggal Nabi yang mengubah penyaluran sentralistis-Madinah menjadi pola kedaerahan menuju Syam, Irak, Mesir, dan Persia.\footnote{Tanjung, "The Role of Tadwin Hadith."; Liew, "The Caliphate Will Last for Thirty Years.".} Para sahabat menetap di kota-kota baru sebagai pengajar, hakim, dan gubernur, sehingga tiap kawasan menumbuhkan khazanah dan kecenderungan pemahaman yang khas.\footnote{Liew, "The Caliphate Will Last for Thirty Years."; Mosa, "Synergizing Structure and Semantics.".} Implikasinya, lahir pewarisan bercorak regional berikut cikal \textit{rihlah}: para pemburu hadis menempuh perjalanan menjumpai sahabat tertentu demi riwayat yang otoritatif.\footnote{Tanjung, "The Role of Tadwin Hadith."; Razak, "Evolution of Hadith Textual Criticism.".} Kedekatan pribadi dan martabat \textit{suhbah} berubah menjadi modal kewibawaan pokok, sebab hadis belum dihimpun dalam kitab besar.\footnote{Musadad, "Articulating the Prophetic Authority."; Mosa, "Synergizing Structure and Semantics.".} Kian beragamnya penyaluran memupuk tuntutan tata cara pembuktian yang berkesadaran metode.\footnote{Razak, "Evolution of Hadith Textual Criticism."; Tanjung, "The Role of Tadwin Hadith.".} Penyebaran tersebut memadat pada empat titik utama --- Madinah yang konservatif sebagai pewaris amalan Nabi, Makkah yang menonjol lewat kewibawaan Ibnu Abbas, serta Kufah dan Basrah sebagai simpul gesit tempat riwayat dan persoalan hukum perkotaan bertemu, hingga perdebatan \textit{ra'yu} kelak memanas.\footnote{Tanjung, "The Role of Tadwin Hadith."; Mosa, "Synergizing Structure and Semantics."; Musadad, "Articulating the Prophetic Authority."; Liew, "The Caliphate Will Last for Thirty Years."; Liew, "The Caliphate Will Last for Thirty Years."; Razak, "Evolution of Hadith Textual Criticism.".} Keterhubungan antarpusat ditopang perjalanan ilmu yang kian gencar.\footnote{Tanjung, "The Role of Tadwin Hadith."; Mosa, "Synergizing Structure and Semantics.".} Titik-titik kewilayahan ini menjelma warisan struktural yang segera diwarisi tabi'in, manakala jalinan personal yang mulai renggang beralih menjadi kelembagaan guru-murid dengan kaidah \textit{isnad} yang kian mengeras.\footnote{Musadad, "Articulating the Prophetic Authority."; Razak, "Evolution of Hadith Textual Criticism.".}

Tuntutan pembuktian dijawab lewat \textit{tathabbut} para Khulafaur Rasyidin dengan ragam prosedur yang bermuara pada tujuan tunggal: menjaga keutuhan informasi.\footnote{Tanjung, "The Role of Tadwin Hadith."; Razak, "Evolution of Hadith Textual Criticism.".} Abu Bakar mensyaratkan kesaksian penyerta (\textit{bayyinah}) sebelum menerima riwayat asing, sebagai wujud tanggung jawab merawat kemurnian ajaran sepeninggal Nabi.\footnote{Razak, "Evolution of Hadith Textual Criticism."; Liew, "The Caliphate Will Last for Thirty Years.".} Umar bahkan menolak dan memperingatkan pembawa hadis yang janggal tanpa penguat, sehingga tercatat sebagai tonggak skeptisisme metodologis terhadap pernyataan perseorangan.\footnote{Tanjung, "The Role of Tadwin Hadith."; Mosa, "Synergizing Structure and Semantics.".} Ali menempuh cara penyumpahan (\textit{yamin}): perawi yang belum dikenal diminta bersumpah atas nama Allah sebagai ujian kejujuran berdimensi psikologis-moral.\footnote{Razak, "Evolution of Hadith Textual Criticism."; Tanjung, "The Role of Tadwin Hadith.".} Rangkaian praktik itu secara tegas merupakan cikal kritik \textit{sanad} --- nalar kritis angkatan pertama atas keterpercayaan sumber, meski belum dirumuskan sebagai disiplin yang terstruktur.\footnote{Mosa, "Synergizing Structure and Semantics."; Razak, "Evolution of Hadith Textual Criticism.".}

Faktor pengakselerasi pelembagaan pembuktian ialah \textit{al-fitnah al-kubra} pascaterbunuhnya Usman, yang memecah umat ke dalam kubu politik-teologis yang saling berebut pengesahan.\footnote{Liew, "The Caliphate Will Last for Thirty Years."; Tanjung, "The Role of Tadwin Hadith.".} Perseteruan Ali--Muawiyah dan kelompok sempalan garis keras menyulut pemalsuan (\textit{wad'}): kisah direkayasa lalu disandarkan kepada Nabi demi menopang klaim sektarian.\footnote{Liew, "The Caliphate Will Last for Thirty Years."; Razak, "Evolution of Hadith Textual Criticism.".} \textit{Afterlife} hadis khilafah tiga puluh tahun membuktikan bahwa satu \textit{matan} dapat diproduksi ulang untuk menopang kisah kekuasaan yang saling bertolak belakang, sehingga krisis keaslian berwatak struktural.\footnote{Liew, "The Caliphate Will Last for Thirty Years."; Musadad, "Articulating the Prophetic Authority.".} Sebagai benteng pertahanan, sahabat dan tabi'in perdana mulai menuntut penyebutan mata rantai penyaluran --- peralihan dari keyakinan mutlak menuju keraguan metodologis.\footnote{Tanjung, "The Role of Tadwin Hadith."; Mosa, "Synergizing Structure and Semantics.".} Jejaknya terekam dalam pernyataan Ibnu Sirin: pemeriksaan \textit{sanad} diberlakukan sesudah fitnah, lengkap dengan kaidah penerimaan berbasis afiliasi --- diambil dari Ahlussunnah dan ditampik dari ahli bidah.\footnote{Liew, "The Caliphate Will Last for Thirty Years."; Razak, "Evolution of Hadith Textual Criticism.".}

\phantomsection
\subsubsection*{3. Periodisasi Masa Tabi'in (80--130 H / 700--750 M)}
\addcontentsline{toc}{subsubsection}{3. Periodisasi Masa Tabi'in (80--130 H / 700--750 M)}

Watak era tabi'in ialah peralihan dari regionalisasi menuju jejaring keilmuan lintas kawasan yang direkatkan oleh \textit{rihlah fi talab al-'ilm}.\footnote{Tanjung, "The Role of Tadwin Hadith."; Mosa, "Synergizing Structure and Semantics.".} Kewibawaan berpindah dari martabat \textit{suhbah} menuju keterandalan jalinan penyaluran guru-ke-murid lintas kota.\footnote{Musadad, "Articulating the Prophetic Authority."; Razak, "Evolution of Hadith Textual Criticism.".} Hijaz menonjol dalam memelihara amalan Madinah, sedangkan Irak tampil dalam wacana hukum-teologi yang tanggap terhadap persoalan perkotaan.\footnote{Tanjung, "The Role of Tadwin Hadith."; Liew, "The Caliphate Will Last for Thirty Years.".} Kegesitan para pencari ilmu meluaskan penyaluran sekaligus mempertemukan ragam riwayat lintas jalur dalam satu forum.\footnote{Mosa, "Synergizing Structure and Semantics."; Razak, "Evolution of Hadith Textual Criticism.".} Era tabi'in dengan demikian menandai tahapan pelembagaan: hadis beranjak dari pusaka personal-komunal menjadi kegiatan keilmuan yang terorganisasi.\footnote{Musadad, "Articulating the Prophetic Authority."; Tanjung, "The Role of Tadwin Hadith.".} Pelembagaan ini disangga tiga wadah: \textit{halaqah} masjid, majelis kediaman ulama, dan jalinan guru-murid melalui \textit{sama'}, \textit{qira'ah}, \textit{ijazah}.\footnote{Tanjung, "The Role of Tadwin Hadith."; Musadad, "Articulating the Prophetic Authority.".} Dalam \textit{halaqah}, syekh membacakan riwayat, murid menyimak, mencatat, lalu mencocokkannya lewat pembacaan ulang demi pengabsahan.\footnote{Musadad, "Articulating the Prophetic Authority."; Razak, "Evolution of Hadith Textual Criticism.".} Majelis khusus memungkinkan penyaluran yang lebih intensif dan selektif, termasuk \textit{ijazah} sebagai lisensi meriwayatkan.\footnote{Tanjung, "The Role of Tadwin Hadith."; Mosa, "Synergizing Structure and Semantics.".} Kekerabatan lintas angkatan melahirkan silsilah terdokumentasi, sehingga \textit{sanad} mencerminkan ikatan pedagogis nyata yang terverifikasi secara historis.\footnote{Mosa, "Synergizing Structure and Semantics."; Musadad, "Articulating the Prophetic Authority.".} Yang tertulis dan yang lisan bekerja berdampingan: catatan murid tetap memerlukan pengabsahan lisan dari guru demi kewibawaan penuh.\footnote{Musadad, "Articulating the Prophetic Authority."; Tanjung, "The Role of Tadwin Hadith.".}

Seiring melebarnya jejaring, \textit{isnad} mengeras dari tuntutan yang longgar menjadi patokan teknis pembuktian, disertai tunas \textit{'ilm al-rijal}.\footnote{Razak, "Evolution of Hadith Textual Criticism."; Mosa, "Synergizing Structure and Semantics.".} Ulama tabi'in mencatat siapa menerima riwayat dari siapa, di mana perjumpaan itu berlangsung, serta mutu hafalan dan integritas tiap mata rantai sebagai arsip penilaian.\footnote{Mosa, "Synergizing Structure and Semantics."; Tanjung, "The Role of Tadwin Hadith.".} Peneguhan ini tak terlepas dari pertikaian sektarian dan derasnya pernyataan riwayat pasca-fitnah yang menuntut atribusi ketat pemisah Ahlussunnah dari ahli bidah.\footnote{Liew, "The Caliphate Will Last for Thirty Years."; Razak, "Evolution of Hadith Textual Criticism.".} Pembandingan ragam \textit{matan} lintas jalur riwayat dijalankan untuk mengendus penyimpangan redaksional --- cikal kritik \textit{sanad-cum-matan} mutakhir.\footnote{Razak, "Evolution of Hadith Textual Criticism."; A. A. R. Al-Imam, "The influence of modernity on matn criticism: A comparative study of traditionist and reformist approaches in Ibn Hajar's Fath al-Bari and Rashid Rida's Majallat al-Manar on hadiths of natural phenomena," \\textit{Journal of Islamic Thought and Civilization}, vol. 16, no. 1, hlm. 01--17, 2026.} Akibatnya, \textit{isnad} tumbuh sebagai jawaban rangkap, ilmiah sekaligus sosiologis, dalam merawat keutuhan teks di tengah perebutan kewibawaan yang tajam.\footnote{Tanjung, "The Role of Tadwin Hadith."; Mosa, "Synergizing Structure and Semantics.".}

Puncak pelembagaan ialah \textit{tadwin} resmi yang disokong negara pada era Umar bin Abdul Aziz (99--101 H/717--720 M), yang risau atas wafatnya para penghafal dan maraknya riwayat yang lemah.\footnote{Tanjung, "The Role of Tadwin Hadith."; Razak, "Evolution of Hadith Textual Criticism.".} Titah resmi kepada para gubernur --- terutama Abu Bakar bin Hazm di Madinah --- dan kepada al-Zuhri memerintahkan penghimpunan, penyaringan, dan pembukuan hadis yang terserak.\footnote{Tanjung, "The Role of Tadwin Hadith."; Musadad, "Articulating the Prophetic Authority.".} Karya perdana umumnya masih berupa \textit{jami'}/\textit{musannaf} bersahaja yang merangkum hadis \textit{marfu'}, fatwa sahabat (\textit{mauquf}), dan pandangan tabi'in (\textit{maqtu'}) ke dalam bab fikih --- belum terhimpun sebagai kompilasi profetik murni seperti abad ketiga.\footnote{Musadad, "Articulating the Prophetic Authority."; Razak, "Evolution of Hadith Textual Criticism.".} \textit{al-Muwatta'} karya Malik memadukan hadis dan amal Madinah, sehingga menunjukkan gerak kewibawaan amalan setempat menuju perekaman tekstual lintas kawasan.\footnote{Musadad, "Articulating the Prophetic Authority."; Tanjung, "The Role of Tadwin Hadith.".} \textit{Tadwin} tampil sebagai titik balik dari yang oral-personal menuju yang skriptural-komunal, sekaligus fondasi bagi ledakan kanonisasi pada era ulama klasik.\footnote{Tanjung, "The Role of Tadwin Hadith."; Razak, "Evolution of Hadith Textual Criticism.".}

\phantomsection
\subsubsection*{4. Periodisasi Masa Ulama Klasik (130--235 H / 750--850 M)}
\addcontentsline{toc}{subsubsection}{4. Periodisasi Masa Ulama Klasik (130--235 H / 750--850 M)}

Watak era ulama klasik ialah kanonisasi: titik kewibawaan berpindah dari perawi perseorangan menuju kitab-kitab kanonik beserta metode kritik yang mapan.\footnote{Musadad, "Articulating the Prophetic Authority."; Razak, "Evolution of Hadith Textual Criticism.".} Penghimpunan menyeleksi ratusan ribu riwayat yang beredar melalui asas pengabsahan yang tegas dan terdokumentasi.\footnote{Tanjung, "The Role of Tadwin Hadith."; Al-Imam, "The Influence of Modernity on Matn Criticism.".} Ragam penulisan berkembang dari \textit{musannaf} campuran menuju himpunan murni (\textit{mujarrad}) yang memilahkan sabda Nabi dari fatwa sahabat dan opini tabi'in, diawali kitab \textit{musnad} seperti \textit{Musnad} Ahmad yang disusun berdasar nama sahabat.\footnote{Musadad, "Articulating the Prophetic Authority."; Razak, "Evolution of Hadith Textual Criticism.".} Puncaknya, \textit{Kutub al-Sittah} --- Bukhari, Muslim, Abu Dawud, Tirmidzi, Nasa'i, Ibnu Majah --- kukuh sebagai rujukan mutlak Sunni hingga era mutakhir.\footnote{Al-Imam, "The Influence of Modernity on Matn Criticism."; Rozikin, "Al-fara'id Is Half of Knowledge.".} Kewibawaan dilembagakan dalam artefak tekstual terstandar yang dapat digandakan, diajarkan, dan diuji lintas angkatan.\footnote{Musadad, "Articulating the Prophetic Authority."; Tanjung, "The Role of Tadwin Hadith.".}

Para tokoh utama mempersonifikasikan patokan tersebut.\footnote{Musadad, "Articulating the Prophetic Authority."; Razak, "Evolution of Hadith Textual Criticism.".} Al-Bukhari (w. 256 H/870 M) menyaring perawi secara ketat dengan mendayagunakan sumber lisan-tertulis melalui \textit{sama'}, \textit{ijazah}, \textit{wajadah}, \textit{'ard}, sehingga keutuhan dan kecermatan menjadi syarat yang tak dapat ditawar.\footnote{Tanjung, "The Role of Tadwin Hadith."; Mosa, "Synergizing Structure and Semantics.".} Muslim (w. 261 H/875 M) menonjol dalam penyandingan lintas jalur riwayat demi keruntutan redaksi, sedangkan Abu Dawud, Nasa'i, dan Ibnu Majah meluaskan kanon ke hukum praktis melalui \textit{sunan}.\footnote{Razak, "Evolution of Hadith Textual Criticism."; Rozikin, "Al-fara'id Is Half of Knowledge.".} Tirmidzi menyumbangkan taksonomi lewat penamaan \textit{hasan} sebagai kategori tengah antara sahih dan \textit{dhaif}, mengatasi kebuntuan dikotomi biner.\footnote{Al-Imam, "The Influence of Modernity on Matn Criticism."; Razak, "Evolution of Hadith Textual Criticism.".} Polemik \textit{afterlife} hadis khilafah tiga puluh tahun membuktikan bahwa kewibawaan para imam kanonik turut menstabilkan ingatan kolektif di tengah perebutan tafsir politik.\footnote{Liew, "The Caliphate Will Last for Thirty Years."; Musadad, "Articulating the Prophetic Authority.".}

Metodologi penyangganya ialah perpaduan kritik \textit{sanad} dan \textit{matan} melalui \textit{al-jarh wa al-ta'dil}.\footnote{Razak, "Evolution of Hadith Textual Criticism."; Al-Imam, "The Influence of Modernity on Matn Criticism.".} Kritik \textit{sanad} menakar kesinambungan mata rantai, \textit{'adalah}, dan \textit{dabt}; kritik \textit{matan} menakar keselarasan terhadap Al-Qur'an, riwayat yang lebih kukuh, akal, dan fakta kesejarahan.\footnote{Al-Imam, "The Influence of Modernity on Matn Criticism."; Rozikin, "Al-fara'id Is Half of Knowledge.".} Telaah bandingan kritik \textit{matan} dalam \textit{Fath al-Bari} (Ibn Hajar) dan \textit{Majallat al-Manar} (Rasyid Rida) menegaskan bahwa kepedulian terhadap isi berakar pada khazanah klasik, bukan reka baru zaman modern --- polarisasi tradisionis-reformis hanya mempertajam rumusannya pada era modern.\footnote{Al-Imam, "The Influence of Modernity on Matn Criticism."; Razak, "Evolution of Hadith Textual Criticism.".} Penyandingan aneka jalur \textit{sanad} dan ragam \textit{matan} ala Muslim membuktikan bahwa kritik klasik berwatak dialektis, tidak berhenti pada nama perawi.\footnote{Razak, "Evolution of Hadith Textual Criticism."; Mosa, "Synergizing Structure and Semantics.".} Anggapan bahwa ulama klasik semata-mata peduli \textit{sanad} dengan demikian merupakan salah kaprah: keduanya saling melengkapi sejak tahap kematangan.\footnote{Al-Imam, "The Influence of Modernity on Matn Criticism."; Tanjung, "The Role of Tadwin Hadith.".} Terobosan epistemologisnya berupa trikotomi sahih-\textit{hasan}-\textit{dhaif} dengan \textit{Sahihayn} sebagai patokan emas.\footnote{Razak, "Evolution of Hadith Textual Criticism."; Al-Imam, "The Influence of Modernity on Matn Criticism.".} Prasyarat penerimaan --- kesinambungan \textit{sanad}, keutuhan dan kecermatan perawi, serta bebas dari \textit{syudzudz} dan \textit{'illah} --- mengalihkan pengabsahan yang intuitif menjadi tata cara intersubjektif yang dapat direplikasi.\footnote{Razak, "Evolution of Hadith Textual Criticism."; Rozikin, "Al-fara'id Is Half of Knowledge.".} Kanonisasi \textit{Sahihayn} turut ditentukan oleh arsitektur kitab yang memudahkan pemakaian dan penerimaan komunitas lintas angkatan.\footnote{Musadad, "Articulating the Prophetic Authority."; Tanjung, "The Role of Tadwin Hadith.".} Daya tahan taksonomi ini terbukti dalam wacana waris mutakhir, yang kesahihan dalilnya masih ditentukan oleh kategori abad ketiga Hijriah.\footnote{Rozikin, "Al-fara'id Is Half of Knowledge."; Al-Imam, "The Influence of Modernity on Matn Criticism.".} Kemapanan itu kemudian membeku menjadi skolastik repetitif, sehingga modernitas --- skeptisisme Barat dan disrupsi digital --- menuntut pembaruan paradigma.\footnote{Liew, "The Caliphate Will Last for Thirty Years."; A. M. B. Matin Bin Salman, "Reconstructing hadith discourse in the digital age: From text to discourse," \\textit{Journal of Ecohumanism}, vol. 4, no. 1, hlm. 1--11, 2025.}

\phantomsection
\subsubsection*{5. Periodisasi Masa Modern (2015--2026)}
\addcontentsline{toc}{subsubsection}{5. Periodisasi Masa Modern (2015--2026)}

Watak era modern, yang terekam dalam literatur 2015--2026, ialah pluralisasi metode yang memadukan kritik \textit{sanad}, \textit{matan}, pembacaan kontekstual, dan ikhtiar digital-komputasional yang interdisipliner.\footnote{Razak, "Evolution of Hadith Textual Criticism."; Matin Bin Salman, "Reconstructing Hadith Discourse in the Digital Age.".} Perkembangan kritik teks memetakan peralihan dari dominasi \textit{sanad} menuju \textit{sanad-cum-matan} yang peka terhadap \textit{asbab al-wurud} dan kemaslahatan.\footnote{Razak, "Evolution of Hadith Textual Criticism."; Al-Imam, "The Influence of Modernity on Matn Criticism.".} Hadis dimengerti tidak sekadar sebagai teks pewarisan, melainkan sebagai wacana yang diproduksi dan dirundingkan dalam ekosistem yang terdesentralisasi.\footnote{Matin Bin Salman, "Reconstructing Hadith Discourse in the Digital Age."; Razak, "Evolution of Hadith Textual Criticism.".} Kritik rangkap atas \textit{sanad}-\textit{matan} menentukan kesahihan dalil legislasi mutakhir, misalnya dalam fikih waris.\footnote{Rozikin, "Al-fara'id Is Half of Knowledge."; Al-Imam, "The Influence of Modernity on Matn Criticism.".} Era modern meluaskan khazanah klasik: etos \textit{tathabbut}, \textit{rihlah}, dan \textit{jarh-ta'dil} dialihbahasakan menjadi perangkat hermeneutik-komputasional baru.\footnote{Razak, "Evolution of Hadith Textual Criticism."; Mosa, "Synergizing Structure and Semantics.".}

Rintangan bersumber pada dua gelombang: skeptisisme orientalis dan guncangan digital.\footnote{Al-Imam, "The Influence of Modernity on Matn Criticism."; Razak, "Evolution of Hadith Textual Criticism.".} Gelombang pertama --- Goldziher dan Schacht melalui teori \textit{projecting back}, yang dimatangkan Juynboll lewat \textit{common link} yang mencurigai ruas di atas periwayat poros\footnote{Razak, "Evolution of Hadith Textual Criticism."; Mosa, "Synergizing Structure and Semantics.".} --- dijawab dengan pembelaan empiris atas kesinambungan \textit{tadwin} abad pertama serta \textit{isnad-cum-matan} yang mengaitkan percabangan \textit{isnad} dengan ragam \textit{matan} lintas geografis.\footnote{Razak, "Evolution of Hadith Textual Criticism."; Al-Imam, "The Influence of Modernity on Matn Criticism.".} Perbandingan tradisionis (\textit{Fath al-Bari}) versus reformis (\textit{Majallat al-Manar}) mengenai hadis gejala alam memperlihatkan polarisasi internal: mempertahankan kanon berhadapan dengan penafsiran ulang yang rasional.\footnote{Al-Imam, "The Influence of Modernity on Matn Criticism."; Razak, "Evolution of Hadith Textual Criticism.".} Perdebatan kaum skeptis dan apologetis justru memacu perumusan kerangka yang lebih canggih ketimbang skolastik yang beku.\footnote{Razak, "Evolution of Hadith Textual Criticism."; Matin Bin Salman, "Reconstructing Hadith Discourse in the Digital Age.".} Gelombang kedua --- mediatisasi lewat media sosial dan AI yang menggeser \textit{sanad} otoritatif menuju kewenangan algoritmik\footnote{Matin Bin Salman, "Reconstructing Hadith Discourse in the Digital Age."; J. Alghamdi et al., "Pretrained models against traditional machine learning for detecting fake hadith," \\textit{Electronics}, vol. 14, no. 17, hlm. 3484, 2025.} --- menghadirkan peredaran infografis dan video ringkas tanpa \textit{takhrij}, sehingga viralitas tanpa keabsahan menggoyang kewibawaan tradisional.\footnote{Matin Bin Salman, "Reconstructing Hadith Discourse in the Digital Age."; Razak, "Evolution of Hadith Textual Criticism.".} Jawabannya bersifat komputasional: model \textit{pretrained} pendeteksi hadis palsu dengan ketepatan bermakna, berpijak pada ciri kebahasaan-struktural\footnote{Alghamdi, "Pretrained Models Against Traditional Machine."; Mosa, "Synergizing Structure and Semantics.".}, ditambah graf pengetahuan-\textit{transformer} untuk disambiguasi narator dan pengendus anomali \textit{sanad} otomatis dalam skala luas.\footnote{Mosa, "Synergizing Structure and Semantics."; Alghamdi, "Pretrained Models Against Traditional Machine.".} Meski begitu, komputasi tetap memerlukan pengabsahan ulama, sebab mesin probabilistik tidak memahami aspek syariat-kesejarahan secara kognitif.\footnote{Alghamdi, "Pretrained Models Against Traditional Machine."; Razak, "Evolution of Hadith Textual Criticism.".}

Sintesis kajian ini menunjukkan bahwa studi hadis melaju menuju \textit{living hadith} dan hermeneutika digital yang menempatkan teknologi sebagai lanjutan etos klasik.\footnote{Matin Bin Salman, "Reconstructing Hadith Discourse in the Digital Age."; Razak, "Evolution of Hadith Textual Criticism.".} \textit{Living hadith} membedah hadis sebagai nilai hidup dalam laku sosial --- dari penafsiran ulang hadis misoginis hingga teologi sebagai daya tahan psikologis.\footnote{Matin Bin Salman, "Reconstructing Hadith Discourse in the Digital Age."; Rozikin, "Al-fara'id Is Half of Knowledge.".} Hermeneutika digital menuntut literasi rangkap: filologis klasik (\textit{takhrij}, kritik \textit{sanad-matan}) sekaligus verifikasi digital atas sumber siber dan deteksi pemalsuan algoritmik.\footnote{Alghamdi, "Pretrained Models Against Traditional Machine."; Mosa, "Synergizing Structure and Semantics.".} Evolusi kritik teks, artikulasi kewenangan profetik dalam \textit{al-Muwatta'} dan \textit{musannaf}, serta peran \textit{tadwin} menegaskan benang merah ketahanan metodologis: dari \textit{sahifah} hingga pangkalan data digital.\footnote{Musadad, "Articulating the Prophetic Authority."; Tanjung, "The Role of Tadwin Hadith."; Razak, "Evolution of Hadith Textual Criticism.".} Pemetaan lima periode ini menjadi landasan analitis bagi uraian penghafalan dan penghimpunan hadis, dari kelisanan hingga kodifikasi resmi.\footnote{Tanjung, "The Role of Tadwin Hadith."; Razak, "Evolution of Hadith Textual Criticism.".}


\phantomsection
\subsection*{B. Penghafalan dan Penghimpunan Hadis}
\addcontentsline{toc}{subsection}{B. Penghafalan dan Penghimpunan Hadis}

\phantomsection
\subsubsection*{1. Tradisi Menghafal Hadis (\textit{Hifz})}
\addcontentsline{toc}{subsubsection}{1. Tradisi Menghafal Hadis (\textit{Hifz})}

Masyarakat Arab pra-Islam meletakkan legitimasi ilmu pada kelisanan: syair disimak, dituturkan, dan dihafal jauh sebelum aksara menyebar meluas.\footnote{H. Rashwan, "Hadith as oral literature through early Islamic literary criticism," \\textit{Studia Islamica}, vol. 119, no. 1, hlm. 34--69, 2024; H. B. Jaiyeoba and N. M. Osmani, "Hadith preservation: Techniques and contemporary efforts," \\textit{Journal of Fatwa Management and Research}, vol. 29, no. 3, hlm. 31--45, 2024.} Praktik itu diwarisi Islam perdana melalui majelis, khotbah, dan perjumpaan langsung dengan Nabi, sehingga sabda, tindakan, dan ketetapan beliau dijaga terutama lewat \textit{hifz}.\footnote{Jaiyeoba, "Hadith Preservation."; S. Watukila, "Historiografi hadis di masa mutaqaddimin dalam pandangan Mustafa Azami," \\textit{Amsal Al-Qur'an: Jurnal Al-Qur'an dan Hadis}, vol. 1, no. 3, hlm. 221--235, 2024.} Penyampaian hadis pada generasi awal bersifat performatif: lafal disertai isyarat tubuh, intonasi suara, dan konteks didaktis yang menghidupkan teks.\footnote{Rashwan, "Hadith As Oral Literature Through."; R. E. Gil, "Tracing the rihla in buldan books: Ibn al-Faqih's muhaddith resources," \\textit{Hitit Theology Journal}, vol. 23, hlm. 140--154, 2024.} Dokumentasi tertulis sebenarnya sudah dikenal, tetapi otoritasnya tetap tersubordinasi pada budaya simak dan koreksi personal; primasi lisan tidak tergantikan.\footnote{M. Lutfianto et al., "History of hadith writing in the precodification period in the prospective of 'Ajjaj al-Khatib," \\textit{Al-Thiqah: Jurnal Ilmu Keislaman}, vol. 7, no. 2, hlm. 308, 2024; Jaiyeoba, "Hadith Preservation.".} Dengan demikian, \textit{hifz} menjadi fondasi kultural bagi transmisi ilmu agama lintas generasi, melampaui sekadar teknik menghafal individual.\footnote{Rashwan, "Hadith As Oral Literature Through."; Watukila, "Historiografi Hadis Di Masa Mutaqaddimin.".} Tradisi ini ditopang oleh perangkat pedagogis baku --- \textit{sama'}, \textit{qira'ah} atau \textit{'ard}, \textit{tasmi'}, \textit{tikrar}, dan \textit{musalsal} yang ritualistik.\footnote{Jaiyeoba, "Hadith Preservation."; Rashwan, "Hadith As Oral Literature Through.".} Dalam \textit{sama'} murid menyerap materi yang didengarnya lewat repetisi; dalam \textit{qira'ah} murid menyetor bacaannya untuk dikoreksi dan diotorisasi guru sebelum berhak meriwayatkan.\footnote{Jaiyeoba, "Hadith Preservation."; Gil, "Tracing the Rihla in Buldan Books.".} Repetisi \textit{tasmi'} mengonsolidasikan hafalan, sedangkan \textit{musalsal} --- pewarisan adab penyampaian sekaligus isi --- menunjukkan bahwa metode periwayatan turut terekam dalam memori transmisi.\footnote{Rashwan, "Hadith As Oral Literature Through."; Watukila, "Historiografi Hadis Di Masa Mutaqaddimin.".} Verifikasi diperkuat oleh lembar \textit{sahifah}, bukti simak, dan kolasi: hafalan disandingkan dengan tulisan sebelum murid dipandang layak mengantongi \textit{ijazah}.\footnote{Lutfianto, "History of Hadith Writing."; Jaiyeoba, "Hadith Preservation.".} Singkatnya, hafalan pada masa awal merupakan kerja kolektif yang diregulasi oleh jaringan guru--murid dan majelis, bukan olah individual yang soliter.\footnote{Gil, "Tracing the Rihla in Buldan Books."; Rashwan, "Hadith As Oral Literature Through.".}

Kelebihan \textit{hifz} terletak pada daya lentur dan vitalitasnya: hadis tetap lestari sebagai tradisi hidup yang kontekstual-didaktis, terus berdenyut kendati infrastruktur tulis masih minim.\footnote{Rashwan, "Hadith As Oral Literature Through."; Jaiyeoba, "Hadith Preservation.".} Sifat performatifnya memungkinkan adaptasi dialektal melalui \textit{bi al-ma'na} --- kelenturan sosiolinguistik yang masuk akal mengingat audiensnya lintas kabilah.\footnote{Rashwan, "Hadith As Oral Literature Through."; Lutfianto, "History of Hadith Writing.".} Sebaliknya, hafalan juga rentan bervariasi: materi dapat terinterpolasi, susunan lafal tertukar, atau nisbah periwayat keliru di sepanjang rantai transmisi.\footnote{Jaiyeoba, "Hadith Preservation."; Watukila, "Historiografi Hadis Di Masa Mutaqaddimin.".} Ingatan lisan mengandung risiko bias sistematis; tanpa kontrol, hafalan semata tidak menjamin otentisitas jangka panjang.\footnote{Rashwan, "Hadith As Oral Literature Through."; S. M. S. Obeid and F. S. Hussein, "Al-Bukhari's sources," \\textit{Jordan Journal for History and Archaeology}, vol. 17, no. 3, hlm. 44--62, 2023.} Inilah paradoks \textit{hifz}: lincah sebagai medium hidup, namun tetap memerlukan topangan arsip tertulis serta kritik perawi agar validitasnya teruji.\footnote{Jaiyeoba, "Hadith Preservation."; Lutfianto, "History of Hadith Writing.".}

Justru dari kerapuhan itulah \textit{hifz} menyumbang pada mekanisme verifikasi: defisit memori dikompensasi oleh \textit{sanad} sebagai teknologi sosial penjamin mutu hafalan tiap perawi.\footnote{Watukila, "Historiografi Hadis Di Masa Mutaqaddimin."; Jaiyeoba, "Hadith Preservation.".} Jaringan \textit{sanad} memiliki dwifungsi --- validasi dan pemetaan sosial: ia merangkai guru--murid, kota-kota ilmu, dan strata generasi, sekaligus mengaudit otoritas setiap riwayat.\footnote{Gil, "Tracing the Rihla in Buldan Books."; Watukila, "Historiografi Hadis Di Masa Mutaqaddimin.".} Mutu hafalan dikaitkan dengan kritik perawi melalui dwisyarat \textit{'adalah} dan \textit{dabt}; perawi yang lemah ingatannya tersisih setelah disandingkan dengan yang lebih kokoh.\footnote{Obeid, "Al-bukhari's Sources."; S. F. Abbas and M. A. Rawabdeh, "The concept of munkar in al-Dhahabi's critique of al-Hakim's hadith authentication in al-Mustadrak," \\textit{Jurnal Studi Ilmu-Ilmu Al-Qur'an dan Hadis}, vol. 26, no. 2, hlm. 545--568, 2025.} Seleksi hafalan semacam ini terbukti menopang kanonisasi al-Bukhari, yang menyeleksi secara ketat akurasi estafet lisan-tulis.\footnote{Obeid, "Al-bukhari's Sources."; Jaiyeoba, "Hadith Preservation.".} Alhasil, tradisi hafalan bukan rival kritik ilmiah, melainkan rahim kelahiran \textit{al-jarh wa al-ta'dil}; dari basis oral inilah wacana menyambung ke dokumentasi era Nabi.\footnote{Rashwan, "Hadith As Oral Literature Through."; Watukila, "Historiografi Hadis Di Masa Mutaqaddimin.".}

\phantomsection
\subsubsection*{2. Penulisan Hadis pada Masa Nabi (\textit{Kitabah})}
\addcontentsline{toc}{subsubsection}{2. Penulisan Hadis pada Masa Nabi (\textit{Kitabah})}

Komunitas Arab pada era Nabi memiliki keberaksaraan yang terbatas: hanya sebagian kecil sahabat yang melek baca-tulis di tengah mayoritas yang bertutur lisan.\footnote{Lutfianto, "History of Hadith Writing."; Watukila, "Historiografi Hadis Di Masa Mutaqaddimin.".} Tulisan berfungsi sebagai alat mnemonik --- nota ringkas, \textit{sahifah} individual, antologi fragmen --- yang mengikat materi seiring jaringan periwayatan terus berekspansi.\footnote{Lutfianto, "History of Hadith Writing."; Jaiyeoba, "Hadith Preservation.".} Migrasi dari hafalan ke tulisan berlangsung secara bertahap, dengan \textit{qira'ah}, \textit{sama'}, dan sesi simak sebagai penghubungnya.\footnote{Lutfianto, "History of Hadith Writing."; Rashwan, "Hadith As Oral Literature Through.".} Historiografi \textit{mutaqaddimin} versi Mustafa Azami menegaskan bahwa penulisan berlangsung kontinu tanpa henti sejak abad pertama Hijriah.\footnote{Watukila, "Historiografi Hadis Di Masa Mutaqaddimin."; Lutfianto, "History of Hadith Writing.".} Artinya, \textit{kitabah} pada masa kenabian melekat dalam ekologi \textit{hifz}, bukan menjadi antitesisnya: amalan \textit{kitabah} menyangga, memantik ingatan, dan mengawal sirkulasi lisan.\footnote{Jaiyeoba, "Hadith Preservation."; Rashwan, "Hadith As Oral Literature Through.".}

Persoalan intinya terletak pada dua rumpun riwayat yang tampak kontradiktif: larangan menulis selain Al-Qur'an, berhadapan dengan otorisasi eksplisit untuk mencatat sabda Nabi.\footnote{Lutfianto, "History of Hadith Writing."; Watukila, "Historiografi Hadis Di Masa Mutaqaddimin.".} Larangan yang paling valid dinisbahkan kepada Abu Sa'id al-Khudri, sedangkan otorisasi dinisbahkan kepada Abdullah bin Amr, yang mengklaim lisannya tidak pernah mengucapkan selain kebenaran.\footnote{Lutfianto, "History of Hadith Writing."; Jaiyeoba, "Hadith Preservation.".} Al-Khatib al-Baghdadi, dalam \textit{Taqyid al-'Ilm}, mengumpulkan kedua kelompok riwayat tersebut dan menyimpulkan bahwa hanya satu riwayat larangan yang valid, berbanding riwayat pembolehan yang jauh lebih banyak.\footnote{Watukila, "Historiografi Hadis Di Masa Mutaqaddimin."; Lutfianto, "History of Hadith Writing.".} Sebagian kecil orientalis memanfaatkan riwayat larangan itu untuk menyimpulkan bahwa hadis baru dibukukan pada abad ketiga --- kesimpulan yang mengabaikan korpus riwayat otorisasi dan bukti \textit{sahifah}.\footnote{Watukila, "Historiografi Hadis Di Masa Mutaqaddimin."; Rashwan, "Hadith As Oral Literature Through.".} Yang dipertaruhkan bukan ada-tidaknya kontroversi itu sendiri, melainkan bagaimana kedua rumpun riwayat yang tampak berseberangan tersebut dibaca secara kritis, historis, dan metodologis.\footnote{Lutfianto, "History of Hadith Writing."; Jaiyeoba, "Hadith Preservation.".} Rekonsiliasi yang disepakati mayoritas ulama memposisikan larangan tersebut sebagai ketentuan yang bersifat kontekstual dan partikular, bukan veto total atas aktivitas dokumentasi.\footnote{Lutfianto, "History of Hadith Writing."; Watukila, "Historiografi Hadis Di Masa Mutaqaddimin.".} Larangan itu menjawab tiga kekhawatiran konkret: kerancuan antara hadis dan wahyu yang sedang turun, teralihkannya perhatian dari primasi wahyu, dan tercampurnya kedua teks pada satu medium.\footnote{Watukila, "Historiografi Hadis Di Masa Mutaqaddimin."; Jaiyeoba, "Hadith Preservation.".} Kuncinya bersifat kronologis: larangan berlaku pada fase awal wahyu, ketika ayat-ayat yang ringkas masih sulit dibedakan dari hadis, sedangkan otorisasi terbit pada fase akhir, ketika sebagian besar wahyu telah turun dan umat mampu mengenali idiom masing-masing.\footnote{Lutfianto, "History of Hadith Writing."; Watukila, "Historiografi Hadis Di Masa Mutaqaddimin.".} Kaidahnya: sirnanya \textit{'illat} meruntuhkan interdiksi, sehingga norma general yang \textit{mujmal} dinasakh oleh otorisasi partikular bagi sahabat yang presisi.\footnote{Watukila, "Historiografi Hadis Di Masa Mutaqaddimin."; Lutfianto, "History of Hadith Writing.".} Begitu preservasi Al-Qur'an mencapai titik optimal, \textit{ijma'} melegitimasi penulisan dan kodifikasi hadis secara terbuka dan luas tanpa keraguan.\footnote{Jaiyeoba, "Hadith Preservation."; Watukila, "Historiografi Hadis Di Masa Mutaqaddimin.".}

Bukti-bukti awal dari berbagai sumber memperkuat rekonsiliasi ini; yang terkuat adalah \textit{Sahifah al-Sadiqah} milik Abdullah bin Amr, yang ditulis atas izin tegas dari Nabi.\footnote{Lutfianto, "History of Hadith Writing."; Watukila, "Historiografi Hadis Di Masa Mutaqaddimin.".} Korpus serupa milik Ali, Jabir, dan sejumlah sahabat lain juga tercatat, disertai arsip diplomasi --- Piagam Madinah dan korespondensi kepada raja-raja asing --- yang menjadi saksi keberaksaraan komunitas mula-mula.\footnote{Watukila, "Historiografi Hadis Di Masa Mutaqaddimin."; Jaiyeoba, "Hadith Preservation.".} Sahabat tidak berhenti pada menghafal: mereka juga membukukan, mempekerjakan penulis, menghadirkan saksi, dan mencatatkan tarikh --- profil penghafal ternyata tidak identik dengan keengganan terhadap aksara.\footnote{Lutfianto, "History of Hadith Writing."; Rashwan, "Hadith As Oral Literature Through.".} Fragmen-fragmen itu beredar lintas generasi sebagai bahan primer bagi antologi kanonik, dengan mata rantai tekstual yang tetap dapat dilacak.\footnote{Obeid, "Al-bukhari's Sources."; Watukila, "Historiografi Hadis Di Masa Mutaqaddimin.".} Arsip awal ini menyediakan basis empiris bagi kajian penghimpunan pada generasi sahabat, saat nota-nota personal ikut berpindah bersama pemiliknya.\footnote{Gil, "Tracing the Rihla in Buldan Books."; Lutfianto, "History of Hadith Writing.".}

\phantomsection
\subsubsection*{3. Penghimpunan Hadis pada Masa Sahabat}
\addcontentsline{toc}{subsubsection}{3. Penghimpunan Hadis pada Masa Sahabat}

Determinan utama penghimpunan pada era sahabat adalah eksodus ke luar Madinah: hafalan dan tulisan turut bermigrasi, melahirkan himpunan-himpunan kewilayahan yang bersifat desentralistis.\footnote{Gil, "Tracing the Rihla in Buldan Books."; Watukila, "Historiografi Hadis Di Masa Mutaqaddimin.".} Ekspansi ke Syam, Irak, Mesir, dan Persia menumbuhkan korpus lokal yang bertumpu pada figur sahabat yang bermukim sebagai pengajar, hakim, maupun gubernur --- ikatan yang bersifat personal dan renggang, bukan sentralisasi pembukuan.\footnote{Gil, "Tracing the Rihla in Buldan Books."; S. Kara, "Debating the origins: The sanctity of Madina in hadith narratives," \\textit{Journal of Islamic Studies}, vol. 37, no. 2, hlm. 163--207, 2026.} Genre \textit{buldan} mendokumentasikan \textit{Rihlah}, yaitu mobilitas penuntut ilmu antarkota menuju sahabat-sahabat kunci --- prototipe perjalanan keilmuan sekaligus agregasi transregional.\footnote{Gil, "Tracing the Rihla in Buldan Books."; Jaiyeoba, "Hadith Preservation.".} Karena energi komunitas lebih tersedot pada pembukuan Al-Qur'an, hadis terakumulasi melalui \textit{sahifah} individual, diktat murid, dan hafalan komunal tanpa patronase negara.\footnote{Lutfianto, "History of Hadith Writing."; Watukila, "Historiografi Hadis Di Masa Mutaqaddimin.".} Dispersi geografis ini memicu divergensi periwayatan, sehingga filtrasi yang ketat menjadi keharusan untuk memurnikan ajaran dari anomali.\footnote{Jaiyeoba, "Hadith Preservation."; Gil, "Tracing the Rihla in Buldan Books.".} Tabi'in menerima estafet berupa nota personal dan memori kolektif sebagai bahan mentah yang kemudian diteruskan, dibandingkan, dan ditata secara sistematis.\footnote{Gil, "Tracing the Rihla in Buldan Books."; Lutfianto, "History of Hadith Writing.".} Kanal utamanya adalah kedekatan guru--murid antara sahabat dan tabi'in, sebagaimana termaktub dalam karya-karya proto seperti \textit{atraf} yang menjadi bekal sebelum majelis.\footnote{Gil, "Tracing the Rihla in Buldan Books."; Jaiyeoba, "Hadith Preservation.".} Prosopografi lingkaran Basrah membuktikan bahwa transmisi pada masa transisi ini lebih difasilitasi oleh tabi'in yang solid ketimbang sahabat lokal, dengan \textit{rihlah} mikro antarsentra yang memperkaya simpul-simpulnya.\footnote{Gil, "Tracing the Rihla in Buldan Books."; Watukila, "Historiografi Hadis Di Masa Mutaqaddimin.".} Lembaran \textit{sahifah}, bukti simak, dan fragmen sahabat dikurasi dalam \textit{halaqah} tabi'in menuju tekstualisasi yang bersifat transregional.\footnote{Lutfianto, "History of Hadith Writing."; Jaiyeoba, "Hadith Preservation.".} Estafet berlapis ini membuka jalan bagi kodifikasi resmi, tatkala penguasa era Umar bin Abdul Aziz menghimpun materi yang berserak ke dalam program yang lebih terencana.\footnote{Watukila, "Historiografi Hadis Di Masa Mutaqaddimin."; Lutfianto, "History of Hadith Writing.".}

Filtrasi ini dikenal dengan istilah \textit{tathabbut}, yaitu etos kehati-hatian Khulafaur Rasyidin yang mensyaratkan bukti tambahan sebelum mengadopsi kabar dari luar.\footnote{Watukila, "Historiografi Hadis Di Masa Mutaqaddimin."; Jaiyeoba, "Hadith Preservation.".} Abu Bakar mensyaratkan \textit{bayyinah} berupa kesaksian tambahan; Umar mendiskualifikasi kabar ganjil yang tidak didukung penguat, bahkan hingga taraf ancaman; sementara Ali mensyaratkan \textit{yamin} bagi perawi yang tidak dikenal.\footnote{Watukila, "Historiografi Hadis Di Masa Mutaqaddimin."; Lutfianto, "History of Hadith Writing.".} Kebijakan diskualifikasi tersebut merintis skeptisisme yang metodis, karena testimoni tunggal mesti lolos uji intersubjektif sebelum dapat diterima sebagai hujah yang sah.\footnote{Jaiyeoba, "Hadith Preservation."; Abbas, "The Concept of Munkar.".} Perluasannya dapat dipetakan pada doktrin \textit{munkar} dalam kritik al-Dzahabi terhadap al-Hakim di \textit{al-Mustadrak}: riwayat yang berlawanan dengan perawi yang lebih superior gugur, sekalipun rantai eksternalnya tetap kontinu.\footnote{Abbas, "The Concept of Munkar."; Obeid, "Al-bukhari's Sources.".} Dengan demikian, \textit{tathabbut} merupakan prototipe kritik \textit{sanad} yang definitif, kendati idiom teknisnya belum baku pada masa itu.\footnote{Watukila, "Historiografi Hadis Di Masa Mutaqaddimin."; Gil, "Tracing the Rihla in Buldan Books.".}

Sahabat senior berfungsi sentral sebagai simpul yang menyambungkan memori kenabian dengan generasi murid tabi'in yang datang kemudian.\footnote{Gil, "Tracing the Rihla in Buldan Books."; Kara, "Debating the Origins.".} Aisyah tampil sebagai otoritas korektif-hermeneutis, yang menyanggah, mengklarifikasi, dan memaknai ulang kesaksian koleganya dengan modal pengalaman langsung dan konsistensi makna; artinya, kritik pemahaman internal telah berdenyut sejak generasi awal.\footnote{Kara, "Debating the Origins."; N. Ardig et al., "Genre analysis and religious texts: A methodological model of hadith commentary," \\textit{Journal of Islamic Studies}, vol. 36, no. 2, hlm. 163--218, 2025.} Ibnu Abbas unggul di Makkah dalam bidang tafsir-fikih, Anas memiliki memori khidmah yang ekstensif, dan Abu Hurairah tampil paling prolifik dengan korpus terluas, meskipun arsitektur kuantitatif jaringannya masih menuntut kajian prosopografi lanjutan.\footnote{Gil, "Tracing the Rihla in Buldan Books."; Watukila, "Historiografi Hadis Di Masa Mutaqaddimin.".} Kontroversi seputar sakralitas Madinah membuktikan bahwa otoritas senior mengonsolidasi memori kolektif di tengah rivalitas tafsir antarregion yang tajam.\footnote{Kara, "Debating the Origins."; Gil, "Tracing the Rihla in Buldan Books.".} Glosa dan koreksi mereka menjadi prototipe syarah yang pada periode berikutnya berkembang menjadi genre otonom dengan tradisi tersendiri.\footnote{Ardig, "Genre Analysis and Religious Texts."; Kara, "Debating the Origins.".}

\phantomsection
\subsubsection*{4. Kodifikasi Hadis Secara Resmi (\textit{Tadwin})}
\addcontentsline{toc}{subsubsection}{4. Kodifikasi Hadis Secara Resmi (\textit{Tadwin})}

Kelahiran \textit{tadwin} dipicu oleh dua kondisi darurat sekaligus: kekhawatiran punahnya khazanah setelah gugurnya \textit{huffaz} di medan perang, dan kontaminasi hadis apokrif-politis pasca-\textit{fitnah}.\footnote{Lutfianto, "History of Hadith Writing."; Watukila, "Historiografi Hadis Di Masa Mutaqaddimin.".} Menjelang akhir abad pertama, para saksi langsung mulai memudar seiring bermekarannya klaim-klaim sektarian; preservasi melalui memori privat dan arsip yang tercecer pun terbukti kurang memadai.\footnote{Watukila, "Historiografi Hadis Di Masa Mutaqaddimin."; Kara, "Debating the Origins.".} Pada era sahabat, pembukuan belum menjadi tuntutan karena fokus masih tersedot pada Al-Qur'an; barulah pada angkatan tabi'in dokumentasi menjadi mendesak guna membendung fabrikasi.\footnote{Lutfianto, "History of Hadith Writing."; Jaiyeoba, "Hadith Preservation.".} Inilah momen infleksi: verifikasi privat-informal model \textit{tathabbut} harus dieskalasi menjadi agenda kolektif yang terorganisasi.\footnote{Jaiyeoba, "Hadith Preservation."; Lutfianto, "History of Hadith Writing.".} Kondisi darurat tersebut memicu munculnya respons negara sebagai solusi institusional atas krisis preservasi yang cakupannya kian meluas.\footnote{Watukila, "Historiografi Hadis Di Masa Mutaqaddimin."; Obeid, "Al-bukhari's Sources.".} Titah Umar bin Abdul Aziz (99--101 H/717--720 M) dialamatkan kepada para gubernur --- terutama Abu Bakar bin Hazm di Madinah --- serta kepada al-Zuhri, dengan perintah untuk mengumpulkan, menyeleksi, dan membukukan hadis yang tercecer.\footnote{Lutfianto, "History of Hadith Writing."; Watukila, "Historiografi Hadis Di Masa Mutaqaddimin.".} Surat khalifah tersebut mengamanatkan penyisiran sistematis atas hafalan dan diktat ulama lintas region --- sebuah intervensi berpatronase negara yang menggeser inisiatif privat menjadi proyek komunal.\footnote{Watukila, "Historiografi Hadis Di Masa Mutaqaddimin."; Jaiyeoba, "Hadith Preservation.".} Al-Zuhri, sebagai eksekutor, mengagregasi bahan antarkota dan merumuskan fondasi metodologisnya, meski ia tidak layak direduksi menjadi poros tunggal kodifikasi dalam episode panjang yang mendahuluinya.\footnote{Lutfianto, "History of Hadith Writing."; Gil, "Tracing the Rihla in Buldan Books.".} \textit{Rihlah} dioptimalkan sebagai wahana ekstraksi, dengan delegasi yang menyambangi repositori hidup di sentra-sentra ilmu yang berpencar.\footnote{Gil, "Tracing the Rihla in Buldan Books."; Lutfianto, "History of Hadith Writing.".} Titah tersebut menjadi prasasti institusional yang menggeser \textit{Tadwin} dari sekadar diskursus menjadi kebijakan negara, sekaligus mematenkan \textit{tadwin} sebagai agenda komunal.\footnote{Watukila, "Historiografi Hadis Di Masa Mutaqaddimin."; Jaiyeoba, "Hadith Preservation.".}

Produk abad kedua masih bersifat transisional: antologi \textit{jami'} dan \textit{musannaf} berformat bab-bab fikih, belum sepenuhnya steril sebagai kumpulan sabda Nabi semata.\footnote{Lutfianto, "History of Hadith Writing."; Obeid, "Al-bukhari's Sources.".} Ketiga kategori --- \textit{marfu'}, \textit{mauquf}, dan \textit{maqtu'} --- masih terbaur, sehingga otoritas profetik belum terdemarkasi dari fikih angkatan awal yang mengapitnya.\footnote{Obeid, "Al-bukhari's Sources."; Ardig, "Genre Analysis and Religious Texts.".} \textit{al-Muwatta'} karya Malik diproyeksikan sebagai manifes konsensus amal Madinah, sementara \textit{Musannaf} karya 'Abd al-Razzaq berekspansi menuju agregasi yang lebih komprehensif, tekstual, dan transregional.\footnote{Obeid, "Al-bukhari's Sources."; Lutfianto, "History of Hadith Writing.".} Korpus transisional ini memasok bahan mentah bagi kodifikasi kanonik abad ketiga: al-Bukhari mendayagunakan sumber lisan-tulis melalui \textit{sama'}, \textit{ijazah}, \textit{wajadah}, dan \textit{'ard}.\footnote{Obeid, "Al-bukhari's Sources."; Abbas, "The Concept of Munkar.".} Amalgamasi kategorial ini, secara ironis, justru membidani kebutuhan akan genre taksonomis baru yang memisahkan graduasi otentisitas secara eksplisit.\footnote{Ardig, "Genre Analysis and Religious Texts."; Obeid, "Al-bukhari's Sources.".}

Konsekuensi historisnya bersifat paradigmatik: supremasi oral-privat bermutasi menjadi rezim skriptural-komunal garapan negara bersama komunitas ilmu.\footnote{Watukila, "Historiografi Hadis Di Masa Mutaqaddimin."; Jaiyeoba, "Hadith Preservation.".} Otoritas bergeser dari kedekatan personal dengan perawi menuju kitab yang direplikasi, disirkulasikan kurikulumnya lintas region, dan diuji melalui kritik \textit{sanad} dan \textit{matan} secara formal.\footnote{Obeid, "Al-bukhari's Sources."; Abbas, "The Concept of Munkar.".} Parameter autentikasi --- kontinuitas \textit{sanad}, integritas dan presisi perawi, serta sterilitas dari \textit{syudzudz} dan \textit{'illah} --- memuncak dalam kanonisasi \textit{Kutub al-Sittah} dan purifikasi \textit{mujarrad}.\footnote{Obeid, "Al-bukhari's Sources."; Jaiyeoba, "Hadith Preservation.".} Kanonisasi ini sekaligus mengonsolidasikan memori kolektif dan mendemarkasi ortodoksi di tengah divergensi transmisi peninggalan \textit{tadwin}.\footnote{Kara, "Debating the Origins."; Ardig, "Genre Analysis and Religious Texts.".} Pascakodifikasi, wacana bermigrasi menuju krisis suksesi: kemapanan arsip tertulis justru menginkubasi fabrikasi mutakhir, yang kemudian menuntut tipologi pemalsuan beserta desain penyelamatannya.\footnote{Abbas, "The Concept of Munkar."; Obeid, "Al-bukhari's Sources.".}


\phantomsection
\subsection*{C. Timbulnya Pemalsuan Hadis dan Upaya Penyelamatan}
\addcontentsline{toc}{subsection}{C. Timbulnya Pemalsuan Hadis dan Upaya Penyelamatan}

\phantomsection
\subsubsection*{1. Pengertian Hadis \textit{Maudu'}}
\addcontentsline{toc}{subsubsection}{1. Pengertian Hadis \textit{Maudu'}}

Kodifikasi formal mengubah transmisi lisan yang tadinya terserak menjadi korpus terdokumentasi, tetapi konsolidasi ini justru membuka celah bagi infiltrasi artifisial sehingga batasan \textit{maudu'} perlu ditegaskan.\footnote{M. A. M. Muhamed Ali et al., "Al-jarh wa al-ta'dil (criticism and praise): Its significance in the science of hadith," \\textit{Mediterranean Journal of Social Sciences}, vol. 6, no. 2S1, hlm. 284, 2015; Wasman et al., "A critical approach to prophetic traditions: Contextual criticism in understanding hadith," \\textit{Al-Jami'ah: Journal of Islamic Studies}, vol. 61, no. 1, hlm. 1--17, 2023.} Hadis \textit{maudu'} dimaknai sebagai pemberitaan fabrikan tanpa basis historis yang dinisbatkan secara dusta kepada Nabi, padahal beliau tidak pernah menyabdakannya.\footnote{Muhamed Ali, "Al-jarh Wa Al-ta'dil (criticism and Praise)."; M. A. Calgan, "Analysis of Ibn Kathir's content criticism of some of the Sahihayn hadiths in the sirah section of al-Bidaya wa al-Nihaya," \\textit{Cumhuriyet Ilahiyat Dergisi}, vol. 28, no. 1, hlm. 303--324, 2024.} Batasan tersebut menegaskan dua unsur pembentuknya: konstruksi \textit{matan} artifisial dan klaim nasab palsu melalui \textit{sanad} rekayasa yang dikesankan tersambung.\footnote{Wasman, "A Critical Approach to Prophetic Traditions."; E. Mohamed and R. Sarwar, "Linguistic features evaluation for hadith authenticity through automatic machine learning," \\textit{Digital Scholarship in the Humanities}, vol. 37, no. 3, hlm. 830--843, 2022.} Berbeda dari kekhilafan yang tak disengaja, \textit{maudu'} mensyaratkan unsur kesengajaan, sehingga penilaiannya mempersoalkan integritas etik transmiter, bukan sekadar daya ingatnya.\footnote{Muhamed Ali, "Al-jarh Wa Al-ta'dil (criticism and Praise)."; M. Shaaban et al., "Prophetic authorial style modeling for detecting fabricated hadiths using AraBERT," \\textit{Journal of Computing \& Biomedical Informatics}, vol. 10, no. 2, 2026.} \textit{Maudu'} karenanya menempati kategori mandiri dalam taksonomi hadis, bukan sekadar tingkatan terendah dari hadis lemah.\footnote{Wasman, "A Critical Approach to Prophetic Traditions."; U. M. Noor, "Muqaddimah of Ibn al-Salah and the revival of hadith studies in the Mamluk era," \\textit{Al-Bayan: Journal of Qur'an and Hadith Studies}, vol. 21, no. 2, hlm. 271--297, 2023.}

Kedudukan \textit{maudu'} menjadi lebih jelas ketika dikontraskan dengan \textit{dhaif}, \textit{matruk}, dan \textit{munkar}, yang masing-masing memuat gradasi kecacatan berbeda.\footnote{Muhamed Ali, "Al-jarh Wa Al-ta'dil (criticism and Praise)."; Wasman, "A Critical Approach to Prophetic Traditions.".} Kategori \textit{dhaif} bersumber dari defek \textit{sanad} --- baik berupa keterputusan rantai maupun rapuhnya hafalan --- meskipun status ini terkadang dapat terangkat melalui penguat lintas jalur atau sokongan yurisprudensial.\footnote{Noor, "Muqaddimah of Ibn Al-salah."; T. Tarmom et al., "Deep learning vs compression-based vs traditional machine learning classifiers to detect hadith authenticity," \\textit{Communications in Computer and Information Science}, hlm. 206--222, 2022.} Sebaliknya, \textit{matruk} dan \textit{munkar} merepresentasikan kerusakan yang lebih serius: transmiternya dituduh berdusta, atau materinya bertentangan dengan riwayat perawi yang jauh lebih superior, sejalan dengan telaah \textit{munkar} atas autentikasi al-Hakim dalam \textit{al-Mustadrak} sebagai indikasi disfungsi \textit{isnad} yang berat.\footnote{Wasman, "A Critical Approach to Prophetic Traditions."; Muhamed Ali, "Al-jarh Wa Al-ta'dil (criticism and Praise).".} Berbeda dari ketiganya, \textit{maudu'} merupakan kreasi deliberatif yang dihimpun tersendiri dalam genre khusus (\textit{kutub al-maudu'at}), sebab terbukti dirancang secara sengaja, bukan sekadar lemah.\footnote{Noor, "Muqaddimah of Ibn Al-salah."; Calgan, "Analysis of Ibn Kathir's Content.".} Pembedaan ini membawa implikasi komputasional: model deteksi otomatis harus mampu memisahkan kerapuhan biasa dari signature linguistik hasil rekayasa.\footnote{A. M. Azmi et al., "Computational and natural language processing based studies of hadith literature: A survey," \\textit{Artificial Intelligence Review}, vol. 52, no. 2, hlm. 1369--1414, 2019; Mohamed, "Linguistic Features Evaluation for Hadith.".}

Konsekuensinya bersifat imperatif: konsensus ulama (\textit{ijma'}) mengharamkan secara mutlak transmisi \textit{maudu'} tanpa penjelasan terbuka mengenai statusnya kepada pendengar.\footnote{Muhamed Ali, "Al-jarh Wa Al-ta'dil (criticism and Praise)."; Wasman, "A Critical Approach to Prophetic Traditions.".} Larangan ini berpijak pada ancaman Nabi berstatus \textit{mutawatir} bahwa pendusta atas namanya telah menyiapkan tempat duduknya di neraka --- formulasi yang merefleksikan jawaban reaktif atas praktik fabrikasi yang telah berlangsung.\footnote{Wasman, "A Critical Approach to Prophetic Traditions."; Noor, "Muqaddimah of Ibn Al-salah.".} Telaah kontekstual menegaskan bahwa kedudukan suatu riwayat menentukan legitimasi pemakaiannya sebagai dasar hukum, materi dakwah, maupun landasan keutamaan amal.\footnote{Wasman, "A Critical Approach to Prophetic Traditions."; Calgan, "Analysis of Ibn Kathir's Content.".} Telaah Ibn Kasir atas sebagian materi \textit{Sahihayn} pada segmen \textit{sirah} membuktikan bahwa korpus kanonik pun tetap tunduk pada evaluasi substansi, sehingga toleransi pemakaian \textit{maudu'} demi \textit{fada'il al-a'mal} tetap ditolak oleh mayoritas ulama.\footnote{Calgan, "Analysis of Ibn Kathir's Content."; Noor, "Muqaddimah of Ibn Al-salah.".} Ketegasan definisi, taksonomi, dan implikasi ini menghantarkan pembahasan pada problem historis: momentum kemunculan awal fabrikasi yang terorganisasi.\footnote{Wasman, "A Critical Approach to Prophetic Traditions."; I.-W. Su, "'Ali b. al-Madini (161-234/778-849): A critical reconstruction of his biography and evaluation of his contribution to hadith criticism," \\textit{Journal of Islamic Studies}, vol. 33, no. 1, hlm. 1--34, 2021.} Kriteria identifikasi klasik terkristalisasi sebagai jawaban atas gelombang pemalsuan yang aktual --- telaah transmiter 'Ali b. al-Madini menguat ketika klaim pasca-fitnah melimpah, sehingga menuntut diferensiasi metodis antara riwayat otentik dan fabrikan.\footnote{Su, "'ali B. Al-madini (161-234/778-849)."; Noor, "Muqaddimah of Ibn Al-salah."; Muhamed Ali, "Al-jarh Wa Al-ta'dil (criticism and Praise).".} Uraian kemudian bertransisi dari sisi definisional ke historis: silsilah fabrikasi sejak era kenabian hingga kritik kontemporer.\footnote{Wasman, "A Critical Approach to Prophetic Traditions."; Su, "'ali B. Al-madini (161-234/778-849).".}

\phantomsection
\subsubsection*{2. Sejarah Munculnya Pemalsuan Hadis}
\addcontentsline{toc}{subsubsection}{2. Sejarah Munculnya Pemalsuan Hadis}

Debat mengenai titik awal fabrikasi mempertemukan dua posisi: segelintir kalangan menempatkan embrionya pada masa Nabi masih hidup, sementara mayoritas memosisikannya pasca-fitnah.\footnote{Wasman, "A Critical Approach to Prophetic Traditions."; Su, "'ali B. Al-madini (161-234/778-849).".} Kelompok minoritas ini bersandar pada peringatan Nabi terhadap para pendusta sebagai indikasi sosiologis bahwa fabrikasi duniawi pernah berlangsung, atau tengah dicoba.\footnote{Wasman, "A Critical Approach to Prophetic Traditions."; Muhamed Ali, "Al-jarh Wa Al-ta'dil (criticism and Praise).".} Namun, kajian filologis-historis yang rigor menunjukkan bahwa laporan spesifik tentang fabrikasi di era Nabi mengandung cacat \textit{sanad} fatal serta pelaku yang tidak jelas, sehingga tidak cukup kuat dijadikan evidensi bagi fabrikasi yang sistematis.\footnote{Su, "'ali B. Al-madini (161-234/778-849)."; Noor, "Muqaddimah of Ibn Al-salah.".} Mayoritas menegaskan bahwa ekosistem sejak era Nabi hingga awal masa Umar relatif steril dari fabrikasi massif, berkat kedekatan temporal dengan wahyu, kontrol internal, dan soliditas moral kolektif sahabat.\footnote{Wasman, "A Critical Approach to Prophetic Traditions."; Muhamed Ali, "Al-jarh Wa Al-ta'dil (criticism and Praise).".} Kajian ini mengikuti pandangan mayoritas tersebut: pemalsuan insidental pada era Nabi tidak setara dengan industri terstruktur yang muncul belakangan.\footnote{Su, "'ali B. Al-madini (161-234/778-849)."; Wasman, "A Critical Approach to Prophetic Traditions.".}

Titik demarkasinya ialah peristiwa \textit{al-fitnah al-kubra}: terbunuhnya Usman dan perang saudara antara Ali dan Muawiyah.\footnote{Wasman, "A Critical Approach to Prophetic Traditions."; Su, "'ali B. Al-madini (161-234/778-849).".} Kekosongan otoritas dan krisis keabsahan yang menyertainya memadatkan kelompok-kelompok radikal --- cikal Syiah fanatik pendukung Ali dan Khawarij yang menentang keduanya --- yang membutuhkan fondasi teologis bagi narasi politik-militer mereka.\footnote{Wasman, "A Critical Approach to Prophetic Traditions."; Noor, "Muqaddimah of Ibn Al-salah.".} Ketika legitimasi dari Qur'an maupun riwayat sahih tidak ditemukan, frustrasi doktrinal ini mendorong rekayasa hadis untuk menjustifikasi posisi sendiri sekaligus meruntuhkan otoritas lawan.\footnote{Wasman, "A Critical Approach to Prophetic Traditions."; Muhamed Ali, "Al-jarh Wa Al-ta'dil (criticism and Praise).".} Irak --- khususnya Kufah dan Basrah --- bertransformasi menjadi sentra produksi terbesar, sehingga otoritas semisal Malik mencurigai ekstremitas transmisi dari Irak dan mewajibkan verifikasi bertingkat.\footnote{Su, "'ali B. Al-madini (161-234/778-849)."; Noor, "Muqaddimah of Ibn Al-salah.".} Ali sendiri, di Siffin, memperingatkan pasukannya tentang pencideraan integritas sabda Nabi.\footnote{Wasman, "A Critical Approach to Prophetic Traditions."; Su, "'ali B. Al-madini (161-234/778-849).".}

Memasuki abad kedua hingga ketiga, seiring konsolidasi Umayyah dan transisi menuju Abbasiyah, fabrikasi meluas dari kontestasi elektoral menjadi industri lintas mazhab dan sekte --- Murji'ah, Qadariyah, Jabariyah --- serta persaingan hegemoni antarmazhab hukum.\footnote{Wasman, "A Critical Approach to Prophetic Traditions."; Noor, "Muqaddimah of Ibn Al-salah.".} Produksi yang tak terkendali ini tersalurkan melalui para \textit{qussas}, ambisi popularitas dan keuangan, serta pasar devosi seputar keutamaan amal, mukjizat, hingga eskatologi.\footnote{Wasman, "A Critical Approach to Prophetic Traditions."; Muhamed Ali, "Al-jarh Wa Al-ta'dil (criticism and Praise).".} Untuk meredam lonjakan atribusi pada abad kedua--ketiga ini, standardisasi kritik \textit{'ilal} dan \textit{rijal} dirintis oleh 'Ali b. al-Madini.\footnote{Su, "'ali B. Al-madini (161-234/778-849)."; Muhamed Ali, "Al-jarh Wa Al-ta'dil (criticism and Praise).".} Abad kedua hingga ketiga dengan demikian merepresentasikan transfigurasi pemalsuan dari instrumen faksi menjadi industri kultural yang meresap ke semua lapisan umat Muslim.\footnote{Wasman, "A Critical Approach to Prophetic Traditions."; Noor, "Muqaddimah of Ibn Al-salah.".} Linimasa ini --- dari embrio insidental melalui politik fitnah hingga industri lintas-sekte --- menegaskan bahwa fabrikasi merupakan simptom krisis otoritas, bukan insiden personal.\footnote{Wasman, "A Critical Approach to Prophetic Traditions."; T. Alam and J. Schneider, "Social network analysis of hadith narrators from Sahih Bukhari," \\textit{Proceedings of 2020 7th IEEE International Conference on Behavioural and Social Computing (BESC)}, hlm. 1--6, 2020.} Pemahaman historis tersebut menjadi prasyarat bagi pembedahan motif, sebab determinan pendorongnya hanya dapat terbaca dalam konteks krisis yang menghasilkannya.\footnote{Wasman, "A Critical Approach to Prophetic Traditions."; Su, "'ali B. Al-madini (161-234/778-849).".} Uraian selanjutnya bergeser dari pertanyaan kapan ke mengapa: klasifikasi motif struktural di balik produksi dan sirkulasi riwayat palsu.\footnote{Muhamed Ali, "Al-jarh Wa Al-ta'dil (criticism and Praise)."; Wasman, "A Critical Approach to Prophetic Traditions.".}

\phantomsection
\subsubsection*{3. Faktor-Faktor yang Mendorong Pemalsuan}
\addcontentsline{toc}{subsubsection}{3. Faktor-Faktor yang Mendorong Pemalsuan}

Variabel politik menempati posisi paling determinan: hadis, sebagai otoritas puncak setelah Al-Qur'an, menghimpun legitimasi besar bagi perebutan kekuasaan.\footnote{Wasman, "A Critical Approach to Prophetic Traditions."; Muhamed Ali, "Al-jarh Wa Al-ta'dil (criticism and Praise).".} Faksi-faksi yang memperebutkan keimamahan, khilafah, dan keutamaan khalifah merancang riwayat tandingan guna menopang kedudukan sekaligus mendiskreditkan lawan --- hadis difungsikan sebagai alat perang propaganda.\footnote{Wasman, "A Critical Approach to Prophetic Traditions."; Su, "'ali B. Al-madini (161-234/778-849).".} Semakin pivotal peran hadis sebagai penentu identitas politik-keagamaan, semakin kuat pula dorongan untuk mencetak tandingannya; nalar pasar otoritas inilah yang menerangkan derasnya fabrikasi di Kufah, Basrah, Wasit, dan Baghdad.\footnote{Wasman, "A Critical Approach to Prophetic Traditions."; Alam, "Social Network Analysis of Hadith.".} Analisis jaringan transmiter \textit{Sahih Bukhari} membuktikan dimensi strukturalnya: pola sentralisasi dan percabangan jalur merefleksikan perebutan pengaruh di lingkaran politik-keilmuan.\footnote{Alam, "Social Network Analysis of Hadith."; A. B. Kamran et al., "SemanticHadith: An ontology-driven knowledge graph for the hadith corpus," \\textit{Journal of Web Semantics}, vol. 78, hlm. 100797, 2023.} Fabrikasi politik karenanya merupakan taktik kolektif kubu-kubu yang memperebutkan otoritas normatif, bukan penyimpangan akhlak perseorangan.\footnote{Wasman, "A Critical Approach to Prophetic Traditions."; Muhamed Ali, "Al-jarh Wa Al-ta'dil (criticism and Praise).".}

Variabel teologis-sektarian memperluas arena politik menjadi konflik doktrinal, meliputi isu \textit{qadar}, kehendak bebas, kedudukan pelaku dosa besar, hingga supremasi antarmazhab fikih.\footnote{Wasman, "A Critical Approach to Prophetic Traditions."; Noor, "Muqaddimah of Ibn Al-salah.".} Kontroversi seputar \textit{qadar} dan iradah dinyatakan secara eksplisit turut memupuk rekayasa dan manipulasi periwayatan, sebab tiap kubu membutuhkan sokongan tekstual profetik.\footnote{Wasman, "A Critical Approach to Prophetic Traditions."; Muhamed Ali, "Al-jarh Wa Al-ta'dil (criticism and Praise).".} Rivalitas antara Sunni, Syiah, dan Khawarij beserta turunannya mendorong penciptaan riwayat pemuliaan tokoh, pelaknatan terhadap seteru, dan klaim eskatologis yang mengokohkan identitas.\footnote{Wasman, "A Critical Approach to Prophetic Traditions."; Su, "'ali B. Al-madini (161-234/778-849).".} Telaah kontekstual menegaskan bahwa narasi polemis semacam ini hanya dapat dievaluasi secara proporsional melalui \textit{asbab al-wurud}, bukan melalui literalisme a-historis.\footnote{Wasman, "A Critical Approach to Prophetic Traditions."; Calgan, "Analysis of Ibn Kathir's Content.".} Kebangkitan tradisi Mamluk melalui \textit{Muqaddimah} karya Ibn al-Salah memperlihatkan bahwa kontroversi warisan senantiasa menuntut pembakuan ulang, karena akumulasi sektarian terus diwariskan antargenerasi.\footnote{Noor, "Muqaddimah of Ibn Al-salah."; Wasman, "A Critical Approach to Prophetic Traditions.".} Determinan destruktif-infiltratif melengkapi gambaran ini: unsur \textit{zindiq} beserta kalangan heretik menginfiltrasi komunitas untuk merusak kewibawaan arus utama dari dalam.\footnote{Wasman, "A Critical Approach to Prophetic Traditions."; Muhamed Ali, "Al-jarh Wa Al-ta'dil (criticism and Praise).".} Sebagian fabrikasi dikaitkan dengan lingkungan menyimpang, termasuk kemungkinan unsur \textit{zindiq} yang menyisipkan laporan penggerus kredibilitas syariat, walaupun saluran operasionalnya tidak selalu terdokumentasi secara rinci.\footnote{Wasman, "A Critical Approach to Prophetic Traditions."; Noor, "Muqaddimah of Ibn Al-salah.".} Berbeda dari kalangan sektarian yang masih mengimani kelompoknya sendiri, fabrikator destruktif berambisi mengaburkan Islam secara keseluruhan, sebab muatannya berkontradiksi dengan prinsip \textit{kulli} syariat.\footnote{Wasman, "A Critical Approach to Prophetic Traditions."; Calgan, "Analysis of Ibn Kathir's Content.".} Kritik \textit{matan} Ibn Kasir atas segmen \textit{sirah} dalam \textit{Sahihayn} relevan untuk mendeteksi infiltrasi semacam ini: laporan yang lolos dari seleksi \textit{sanad} tetap dapat digugurkan melalui pengujian koherensi historis-teologis.\footnote{Calgan, "Analysis of Ibn Kathir's Content."; Wasman, "A Critical Approach to Prophetic Traditions.".} Pemodelan gaya profetik berbasis \textit{AraBERT} menyediakan instrumen modern untuk mendeteksi penyimpangan stilistik infiltratif secara statistik.\footnote{Shaaban, "Prophetic Authorial Style Modeling."; Mohamed, "Linguistic Features Evaluation for Hadith.".}

Variabel keempat, yang bercorak sosio-ekonomi-devosional --- mencakup motif dakwah yang keliru melalui \textit{targhib-tarhib}, kepentingan material, dan kebutuhan akan literatur populer --- justru memikul tanggung jawab terbesar atas bertahannya sirkulasi \textit{maudu'} hingga sekarang.\footnote{Wasman, "A Critical Approach to Prophetic Traditions."; Muhamed Ali, "Al-jarh Wa Al-ta'dil (criticism and Praise).".} Fabrikator kerap berdalih konstruktif: laporan palsu tentang keutamaan amal dinilai bermanfaat untuk mendorong kesalehan, sehingga kreasi demi dakwah yang mulia dianggap terlegitimasi.\footnote{Wasman, "A Critical Approach to Prophetic Traditions."; Noor, "Muqaddimah of Ibn Al-salah.".} Para \textit{qussas}, forum penceritaan, dan dakwah populer berperan sebagai saluran distribusi yang efektif dalam memenuhi hasrat moral-emosional-eskatologis di luar supervisi ahli hadis.\footnote{Wasman, "A Critical Approach to Prophetic Traditions."; Muhamed Ali, "Al-jarh Wa Al-ta'dil (criticism and Praise).".} Motivasi material memperburuk keadaan: nama Nabi dimanfaatkan untuk mendulang popularitas dan sokongan finansial dari penguasa maupun audiens.\footnote{Muhamed Ali, "Al-jarh Wa Al-ta'dil (criticism and Praise)."; Wasman, "A Critical Approach to Prophetic Traditions.".} Kombinasi antara motif ideologis di pihak produsen dan motif populer di pihak sirkulator menjelaskan mengapa kemapanan kritik tidak serta-merta memberantas \textit{maudu'} --- pembahasan selanjutnya berlanjut pada ciri-ciri hadis palsu.\footnote{Wasman, "A Critical Approach to Prophetic Traditions."; Azmi, "Computational and Natural Language Processing Based Studies of Hadith Literature.".}

\phantomsection
\subsubsection*{4. Ciri-Ciri Hadis Palsu}
\addcontentsline{toc}{subsubsection}{4. Ciri-Ciri Hadis Palsu}

Kriteria \textit{sanad} bertumpu pada reliabilitas transmiter, kontinuitas mata rantai, dan pendeteksian cacat (\textit{'illah}) melalui disiplin \textit{al-jarh wa al-ta'dil}.\footnote{Muhamed Ali, "Al-jarh Wa Al-ta'dil (criticism and Praise)."; Su, "'ali B. Al-madini (161-234/778-849).".} Indikator terkuatnya meliputi konfesi eksplisit dari pelaku, transmiter yang dikenal pendusta atau pemalsu, pretensi transmisi langsung yang berseberangan dengan catatan biografis, diskontinuitas rantai, dan transmiter \textit{majhul}.\footnote{Muhamed Ali, "Al-jarh Wa Al-ta'dil (criticism and Praise)."; Wasman, "A Critical Approach to Prophetic Traditions.".} Praktik pemeriksaan tanggal lahir dan wafat, inventarisasi hubungan guru-murid, serta verifikasi \textit{liqa'} sejak generasi awal telah dibuktikan oleh 'Ali b. al-Madini, yang menggugurkan klaim perjumpaan yang mustahil secara kronologis.\footnote{Su, "'ali B. Al-madini (161-234/778-849)."; Muhamed Ali, "Al-jarh Wa Al-ta'dil (criticism and Praise).".} Analisis jaringan transmiter \textit{Sahih Bukhari} memodernkan pendekatan ini: anomali relasional --- baik berupa keterisolasian dari gugus seangkatan maupun trayektori yang deviatif dari pola kolektif --- dibaca sebagai isyarat statistik adanya distorsi.\footnote{Alam, "Social Network Analysis of Hadith."; Kamran, "Semantichadith.".} Ontologi \textit{SemanticHadith} memformalkan pendekatan tersebut: transmiter, karya, dan jalur direpresentasikan sebagai jaringan semantik sehingga anomalinya dapat dikalkulasi.\footnote{Kamran, "Semantichadith."; Azmi, "Computational and Natural Language Processing Based Studies of Hadith Literature.".}

Kriteria \textit{matan} menguji substansi riwayat terhadap Al-Qur'an, hadis sahih, fakta sejarah, akal, konsensus ulama, dan prinsip umum syariat.\footnote{Wasman, "A Critical Approach to Prophetic Traditions."; Calgan, "Analysis of Ibn Kathir's Content.".} Gejala kewaspadaannya meliputi: absennya riwayat dalam koleksi yang diakui, redaksi yang janggal atau berbahasa rusak, makna yang mustahil secara ontologis, kontradiksi dengan evidensi yang lebih kuat, konten sektarian atau tidak senonoh, serta imbalan pahala atau ancaman yang tidak proporsional untuk amal yang sepele.\footnote{Wasman, "A Critical Approach to Prophetic Traditions."; Mohamed, "Linguistic Features Evaluation for Hadith.".} Kritik Ibn Kasir atas segmen \textit{sirah} dalam \textit{Sahihayn} membuktikan bahwa kekuatan \textit{sanad} tidak selalu final: rantai yang kukuh pun dapat gugur melalui pengujian koherensi historis.\footnote{Calgan, "Analysis of Ibn Kathir's Content."; Wasman, "A Critical Approach to Prophetic Traditions.".} Evaluasi fitur linguistik melalui \textit{machine learning} menegaskan relevansi kontemporer pendekatan ini: sintaksis, diversitas kosakata, dan deviasi stilistik terbukti berkorelasi dengan otentisitas.\footnote{Mohamed, "Linguistic Features Evaluation for Hadith."; Shaaban, "Prophetic Authorial Style Modeling.".} Pemodelan gaya profetik dengan \textit{AraBERT} memperkuat identifikasi ini melalui sidik stilistik profetik sebagai penakar deviasi.\footnote{Shaaban, "Prophetic Authorial Style Modeling."; F. Haque et al., "Hadith authenticity prediction using sentiment analysis and machine learning," \\textit{2020 IEEE 14th International Conference on Application of Information and Communication Technologies (AICT)}, hlm. 1--6, 2020.}

Pendekatan rekomendasi bukanlah parameter tunggal, melainkan sintesis kumulatif dan kontekstual antara \textit{sanad}, \textit{matan}, dan latar historis.\footnote{Muhamed Ali, "Al-jarh Wa Al-ta'dil (criticism and Praise)."; Wasman, "A Critical Approach to Prophetic Traditions.".} Rantai yang lemah hanya menetapkan status \textit{dhaif}; vonis fabrikasi yang meyakinkan memerlukan konvergensi bukti berupa transmiter fraudulen, kontradiksi tekstual, absennya dukungan historis-sumber, transmisi yang mustahil, dan pengakuan pelaku.\footnote{Wasman, "A Critical Approach to Prophetic Traditions."; Calgan, "Analysis of Ibn Kathir's Content.".} Kritik kontekstual menekankan bahwa \textit{sabab wurud} turut perlu diuji, sebab telaah \textit{sanad}-\textit{matan} saja tidak selalu menyingkap fabrikasi yang dipoles rapi.\footnote{Wasman, "A Critical Approach to Prophetic Traditions."; Noor, "Muqaddimah of Ibn Al-salah.".} Perbandingan klasifier \textit{deep learning}, kompresi, dan \textit{machine learning} tradisional memperlihatkan keunggulan kombinasi fitur struktural-semantik dibanding pendekatan tunggal --- selaras dengan kaidah kumulatif klasik.\footnote{Tarmom, "Deep Learning Vs Compression-based Vs."; Haque, "Hadith Authenticity Prediction Using Sentiment.".} Tinjauan komputasional-NLP mengonfirmasi bahwa prasarana semantik modern meneguhkan nalar klasik: autentikasi menuntut multi-metode bagi setiap laporan.\footnote{Azmi, "Computational and Natural Language Processing Based Studies of Hadith Literature."; Kamran, "Semantichadith.".} Penguasaan ciri \textit{sanad}-\textit{matan} mengharuskan institusionalisasi respons, sebab identifikasi personal tidak mencukupi untuk menghadapi industri yang sistematis.\footnote{Muhamed Ali, "Al-jarh Wa Al-ta'dil (criticism and Praise)."; Noor, "Muqaddimah of Ibn Al-salah.".} Identifikasi otomatis melalui analisis sentimen dan \textit{machine learning} hanya signifikan bila dilatih dengan taksonomi klasik --- pembaruan algoritmik tetap bertumpu pada khazanah kriteria ulama.\footnote{Haque, "Hadith Authenticity Prediction Using Sentiment."; Tarmom, "Deep Learning Vs Compression-based Vs.".} Kebangkitan tradisi Mamluk melalui \textit{Muqaddimah} karya Ibn al-Salah membuktikan bahwa setiap krisis direspons melalui kodifikasi metodologis yang segar, sebuah pola repetitif dari era klasik hingga digital.\footnote{Noor, "Muqaddimah of Ibn Al-salah."; Su, "'ali B. Al-madini (161-234/778-849).".} Telaah Ibn Kasir terhadap periwayatan kanonik memperlihatkan kultur kritik yang tanpa area imun: keterbukaan pada pengujian ulang itulah yang menjadi proteksinya.\footnote{Calgan, "Analysis of Ibn Kathir's Content."; Wasman, "A Critical Approach to Prophetic Traditions.".} Berangkat dari kebutuhan institusionalisasi tersebut, uraian berpindah menuju arsitektur penyelamatan: kerangka metodologis bagi pemurnian korpus secara permanen.\footnote{Muhamed Ali, "Al-jarh Wa Al-ta'dil (criticism and Praise)."; Noor, "Muqaddimah of Ibn Al-salah.".}

\phantomsection
\subsubsection*{5. Upaya Penyelamatan Hadis}
\addcontentsline{toc}{subsubsection}{5. Upaya Penyelamatan Hadis}

Upaya pertama ialah \textit{takhrij}: penelusuran asal genetik sekaligus verifikasi terpadu antara \textit{sanad} dan \textit{matan}.\footnote{Muhamed Ali, "Al-jarh Wa Al-ta'dil (criticism and Praise)."; Wasman, "A Critical Approach to Prophetic Traditions.".} Seluruh jalur suatu \textit{matan} ditelusuri oleh peneliti, varian antarjalur dikomparasikan, kredibilitas tiap transmiter dinilai, dan koherensinya terhadap Qur'an, hadis sahih, serta fakta historis diuji sebelum vonis dijatuhkan.\footnote{Wasman, "A Critical Approach to Prophetic Traditions."; Su, "'ali B. Al-madini (161-234/778-849).".} Tahapan \textit{takhrij} kini dibakukan pula sebagai fitur komputasional --- asesmen linguistik berbasis \textit{machine learning} dan prediksi sentimen melipatgandakan kecepatan penelusuran lintas-korpus tanpa meninggalkan nalar klasik.\footnote{Mohamed, "Linguistic Features Evaluation for Hadith."; Haque, "Hadith Authenticity Prediction Using Sentiment.".} Ontologi \textit{SemanticHadith} mengotomatisasi proses ini: korpus direpresentasikan sebagai jaringan semantik yang tertelusur, terbandingkan, dan teraudit secara sistematis.\footnote{Kamran, "Semantichadith."; Azmi, "Computational and Natural Language Processing Based Studies of Hadith Literature.".} Dengan demikian, \textit{takhrij} berfungsi sebagai prosedur operasional yang menerjemahkan kriteria teoretis tentang ciri \textit{maudu'} menjadi putusan praktis bagi setiap riwayat yang bersirkulasi.\footnote{Muhamed Ali, "Al-jarh Wa Al-ta'dil (criticism and Praise)."; Wasman, "A Critical Approach to Prophetic Traditions.".} Upaya kedua ialah \textit{al-jarh wa al-ta'dil}: penaksiran \textit{'adalah} dan \textit{dabt} setiap mata rantai sebagai benteng epistemologis.\footnote{Muhamed Ali, "Al-jarh Wa Al-ta'dil (criticism and Praise)."; Su, "'ali B. Al-madini (161-234/778-849).".} Disiplin ini semata-mata mengesahkan transmisi dari transmiter \textit{thiqah} sebagai hujah, sembari memfiltrasi sosok \textit{majhul}, pelupa berat, dan pendusta sebelum mencemari korpus normatif.\footnote{Muhamed Ali, "Al-jarh Wa Al-ta'dil (criticism and Praise)."; Noor, "Muqaddimah of Ibn Al-salah.".} Telaah perintis 'Ali b. al-Madini meletakkan dasar bagi \textit{'ilal} dan \textit{rijal} yang kemudian diadopsi oleh Ibn al-Salah dan era Mamluk sebagai kurikulum baku.\footnote{Su, "'ali B. Al-madini (161-234/778-849)."; Noor, "Muqaddimah of Ibn Al-salah.".} Analisis jaringan transmiter \textit{Sahih Bukhari} memvalidasinya secara kuantitatif: transmiter poros yang tertaut padat dinilai lebih andal dibanding simpul pinggiran yang menyimpang.\footnote{Alam, "Social Network Analysis of Hadith."; Kamran, "Semantichadith.".} Melalui \textit{jarh-ta'dil}, komunitas merancang pertahanan bertingkat yang terekam dalam literatur tersendiri.\footnote{Muhamed Ali, "Al-jarh Wa Al-ta'dil (criticism and Praise)."; Wasman, "A Critical Approach to Prophetic Traditions.".}

Upaya ketiga ialah kodifikasi \textit{maudu'at}: penghimpunan laporan-laporan palsu menjadi arsip anti-disinformasi.\footnote{Noor, "Muqaddimah of Ibn Al-salah."; Calgan, "Analysis of Ibn Kathir's Content.".} Kontribusi paling tegas disumbangkan oleh Ibn al-Jauzi melalui \textit{al-Maudhu'at}, yang menghimpun sekitar 1.600 laporan palsu --- sebelumnya terserak di berbagai literatur \textit{rijal} --- menjadi kritik tematik yang sistematis.\footnote{Noor, "Muqaddimah of Ibn Al-salah."; Wasman, "A Critical Approach to Prophetic Traditions.".} Tradisi ini dilanjutkan oleh al-Suyuti (\textit{al-La'ali al-Masnu'ah}), Ibn 'Iraq (\textit{Tanzih al-Syari'ah}), Syaukani (\textit{al-Fawa'id al-Majmu'ah}), dan pada masa modern oleh al-Albani (\textit{Silsilah al-Ahadits al-Dha'ifah wa al-Maudhu'ah}) yang menghidupkan kembali sensibilitas verifikasi revivalis.\footnote{Noor, "Muqaddimah of Ibn Al-salah."; Muhamed Ali, "Al-jarh Wa Al-ta'dil (criticism and Praise).".} Kritik Ibn Kasir yang mengoreksi sebagian materi \textit{Sahihayn} membuktikan bahwa pembakuan kritis mencakup kanon tanpa kecuali.\footnote{Calgan, "Analysis of Ibn Kathir's Content."; Wasman, "A Critical Approach to Prophetic Traditions.".} Korpus \textit{maudu'at} klasik kini memiliki dwifungsi sebagai bahan latih model komputasional --- khazanah klasik menopang pembaruan digital.\footnote{Tarmom, "Deep Learning Vs Compression-based Vs."; Shaaban, "Prophetic Authorial Style Modeling.".}

Sebagai rangkuman Bab II: fabrikasi bukanlah deviasi pinggiran, melainkan krisis teologis kolosal yang mendesak peradaban Islam menyusun pertahanan epistemologis paling mutakhir --- terbentang dari \textit{tathabbut} para sahabat hingga graf pengetahuan digital masa kini.\footnote{Muhamed Ali, "Al-jarh Wa Al-ta'dil (criticism and Praise)."; Kamran, "Semantichadith.".} Survei komputasional, asesmen linguistik otomatis, dan pemodelan \textit{AraBERT} menunjukkan bahwa pembaruan algoritmik berperan sebagai penerus etos \textit{takhrij} dan \textit{jarh-ta'dil}, bukan sebagai substitusi bagi ulama.\footnote{Azmi, "Computational and Natural Language Processing Based Studies of Hadith Literature."; Mohamed, "Linguistic Features Evaluation for Hadith."; Shaaban, "Prophetic Authorial Style Modeling.".} Akan tetapi, digitalisasi turut mengakselerasi peredaran kutipan yang belum terverifikasi, sehingga dependensi terhadap penilaian tradisional \textit{isnad}-\textit{matan} justru semakin menguat.\footnote{Azmi, "Computational and Natural Language Processing Based Studies of Hadith Literature."; Wasman, "A Critical Approach to Prophetic Traditions.".} Setelah periodisasi, penghimpunan, pemalsuan, dan penyelamatan dipetakan, uraian ini siap disintesiskan dalam Bab III: kesimpulan integratif dan agenda verifikasi digital yang berlandaskan pada kritik klasik.\footnote{Noor, "Muqaddimah of Ibn Al-salah."; Kamran, "Semantichadith.".}

